\chapter*{Abstract}
\noindent
Debugging embedded systems requires instruments able to capture digital
signals with tight timing precision, yet the cost of commercial logic
analyzers restricts access to them in educational and research settings.
A low-cost data acquisition (DAQ) device based on the RP2350
microcontroller, compatible with the Sigrok/PulseView ecosystem, offloads
sampling to the Programmable I/O (PIO) subsystem and data movement to the
Direct Memory Access (DMA) controller; however, the benefit of this
architecture over the conventional alternative, polling the general-purpose
I/O (GPIO) pins with the CPU itself, had not been quantified. This work
implements a second firmware version that performs acquisition entirely in
software on a single core, keeping the host communication protocol, the USB stack and the
RLE encoding identical, and experimentally compares both architectures in
terms of maximum sustained sample rate, sampling jitter, sample loss during
continuous capture and CPU load. Measurements used a reference square wave,
a commercial logic analyzer as an external reference and 30 repetitions per
condition; the two architectures were compared using the Mann--Whitney
test with Holm--Bonferroni correction. The PIO+DMA architecture sustained 3.66\,MS/s with an exact
edge count, jitter equal to the quantization limit and no lost samples in
continuous capture up to 3.54\,MS/s, aborting the capture with an explicit
error outside this range, with CPU load between 0.08\,\% and 83\,\%,
proportional to the sample rate. CPU polling occupied the processor entirely
and, although its loop cost would allow about 4.29\,MS/s, it failed
silently from $\approx 1.1$\,MS/s onwards: the loop stalled while the same
CPU transmitted the data, which produced extra edges in triggered captures,
jitter 106 to 400\,ns above quantization and the loss of 24\,\% to 43\,\%
of the samples in continuous capture. Single-core CPU polling, with the same
core both sampling and transmitting, is therefore sufficient for short
captures at rates up to hundreds of kS/s, while PIO+DMA is
required at higher rates, in continuous capture or when the CPU must remain
free. Beyond sample rate, the architecture also sets how costly
failure detection is: under PIO+DMA, queue overflow is a hardware event the
CPU observes outside the critical path, whereas under software acquisition,
detecting the pauses consumes cycles of the sampling loop itself.

\vspace{5mm}
\noindent\textbf{\textit{Keywords}:} Data acquisition. Microcontrollers. RP2350. PIO. DMA.
